\documentclass[UTF8,11pt]{ctexart}
\usepackage[a4paper,margin=22mm]{geometry}
\usepackage{amsmath,amssymb,bm,graphicx,adjustbox,fvextra,longtable,array}
\xeCJKDeclareCharClass{CJK}{"2160 -> "217F, "2460 -> "249B, "2200 -> "22FF}
\newcommand{\fitmath}[1]{\begin{adjustbox}{max width=\linewidth}$\displaystyle #1$\end{adjustbox}}
\setlength{\parindent}{0pt}\setlength{\parskip}{.65em}\renewcommand{\arraystretch}{1.3}
\sloppy\allowdisplaybreaks
\begin{document}
\section*{2015 年考研政治真题}
\subsection*{2015 年全国硕士研究生入学统一考试思想政治理论试题}

\subsection*{一、单项选择题：1～16 小题，每小题1 分，共16 分。下列每题给出的四个选项中，只有一}

个选项是符合题目要求的。请在答题卡上将所选项的字母涂黑。


1.院士袁隆平曾结合自己的科研经历，语重心长地对年轻人说：“书本知识非常重要，电脑技术也很重要，但是书本电脑里种不出水稻来，只有在田里才能种出水稻来。”这表明：


\begin{center}\begin{tabular}{>{\raggedright\arraybackslash}p{\dimexpr(\linewidth-4\tabcolsep)/2\relax}>{\raggedright\arraybackslash}p{\dimexpr(\linewidth-4\tabcolsep)/2\relax}}
\textbf{A.} 实践水平的提高有赖于认识水平的提高 & \textbf{B.} 实践是人类认识的基础和来源\\[.35em]
\textbf{C.} 理论对实践的指导作用没有正确与错误之分 & \textbf{D.} 由实践到认识的第一次飞跃比认识到实践的第二次飞跃更加重要\\[.35em]\end{tabular}\end{center}

2.社会意识，社会意识是社会存在的反映。社会意识具有相对独立性，即它在反映社会存在的同时，还有自己特有的发展形式和规律。社会意识相对独立性最突出的表现是：


\begin{center}\begin{tabular}{>{\raggedright\arraybackslash}p{\dimexpr(\linewidth-4\tabcolsep)/2\relax}>{\raggedright\arraybackslash}p{\dimexpr(\linewidth-4\tabcolsep)/2\relax}}
\textbf{A.} 社会意识内部各种形式之间的相互作用和影响 & \textbf{B.} 社会意识与社会存在发展的不完全同步性\\[.35em]
\textbf{C.} 社会意识各种形式各自具有其历史继承性 & \textbf{D.} 社会意识对社会存在具有能动的反作用\\[.35em]\end{tabular}\end{center}

3.第二次世界大战以后，资本主义国家经历了第三次科技革命，机器大工业发展到自动化阶段。只能化工厂创造出了较高的生产效率，显露出巨大的竞争力。企业在“机器换人”中取得了一定的经济效益。这意味着率先使用机器人的个别企业：


\begin{center}\begin{tabular}{>{\raggedright\arraybackslash}p{\dimexpr(\linewidth-4\tabcolsep)/2\relax}>{\raggedright\arraybackslash}p{\dimexpr(\linewidth-4\tabcolsep)/2\relax}}
\textbf{A.} 剩余价值来源的改变 & \textbf{B.} 资本技术构成的提高\\[.35em]
\textbf{C.} 所生产商品价值提高 & \textbf{D.} 获得更多的社会平均利润\\[.35em]\end{tabular}\end{center}

4.马克思主义政党是科学社会主义与工人运动相结合的产物，是工人阶级的先锋队，这表明：


\begin{center}\begin{tabular}{>{\raggedright\arraybackslash}p{\dimexpr(\linewidth-4\tabcolsep)/2\relax}>{\raggedright\arraybackslash}p{\dimexpr(\linewidth-4\tabcolsep)/2\relax}}
\textbf{A.} 马克思主义政党以工人阶级为基础 & \textbf{B.} 马克思主义政党即工人阶级本身\\[.35em]
\textbf{C.} 马克思主义政党的阶级性是其先进性的根本前提 & \textbf{D.} 马克思主义政党的先进性决定了工人阶级的先进性\\[.35em]\end{tabular}\end{center}

5.新中国的工业化是在苏联的影响下起步的。走中国工业化道路，是中国共产党初步探索我国社会主义建设道路的一个重要思想。当时所讲的工业化道路问题，主要是指：


\begin{center}\begin{tabular}{>{\raggedright\arraybackslash}p{\dimexpr(\linewidth-4\tabcolsep)/2\relax}>{\raggedright\arraybackslash}p{\dimexpr(\linewidth-4\tabcolsep)/2\relax}}
\textbf{A.} 经济建设和国防建设的关系问题 & \textbf{B.} 中央和地方的关系问题\\[.35em]
\textbf{C.} 沿海工业和内地工业的关系问题 & \textbf{D.} 重工业、轻工业和农业的发展关系问题\\[.35em]\end{tabular}\end{center}

6.十三大召开前夕，邓小平强调指出：“社会主义本身是共产主义的初阶段，而我们中国又处在社会主义的初级阶段，就是不发达的阶段。一切都要从这个实际出发，根据这个实际来制定规划。”这一论述：


\begin{center}\begin{tabular}{>{\raggedright\arraybackslash}p{\dimexpr(\linewidth-4\tabcolsep)/2\relax}>{\raggedright\arraybackslash}p{\dimexpr(\linewidth-4\tabcolsep)/2\relax}}
\textbf{A.} 首次提出了社会主义初级阶段概念 & \textbf{B.} 首次系统阐述了社会主义初级阶段理论\\[.35em]
\textbf{C.} 首次把社会主义初级阶段作为事关全局的基本国情加以把握 & \textbf{D.} 首次对社会主义发展阶段进行了划分\\[.35em]\end{tabular}\end{center}

7.城镇化是现代化的必由之路，解决好人的问题是推进城镇化的关键。当前，我国实现城镇化的首要任务是：


\begin{center}\begin{tabular}{>{\raggedright\arraybackslash}p{\dimexpr(\linewidth-4\tabcolsep)/2\relax}>{\raggedright\arraybackslash}p{\dimexpr(\linewidth-4\tabcolsep)/2\relax}}
\textbf{A.} 推进农业转移人口市民化 & \textbf{B.} 使土地的城镇化优于人口的城镇化\\[.35em]
\textbf{C.} 促进农村劳动力向非农产业转移 & \textbf{D.} 实现“人的无差别发展”\\[.35em]\end{tabular}\end{center}

8.1997 年7 月1 日，中国政府对香港恢复行使主权，香港特别行政区成立，香港特别行政区基本法开始实施。香港进入“一国两制”、“港人治港”、高度自治的历史新纪元。香港特别行政区的高度自治权是:


\begin{center}\begin{tabular}{>{\raggedright\arraybackslash}p{\dimexpr(\linewidth-4\tabcolsep)/2\relax}>{\raggedright\arraybackslash}p{\dimexpr(\linewidth-4\tabcolsep)/2\relax}}
\textbf{A.} 特别行政区的完全自治 & \textbf{B.} 中央授权之外的剩余权力\\[.35em]
\textbf{C.} 特别行政区本身固有的权力 & \textbf{D.} 中央授予的地方事务管理权\\[.35em]\end{tabular}\end{center}

9.近代中国，一些爱国人士提出过工业救国、教育救国、科学救国等主张，并为此进行过努力，但这些主张并不能从根本上给濒临危亡的中国指明正确的出路，这是因为他们没有认识到 :


\begin{center}\begin{tabular}{>{\raggedright\arraybackslash}p{\dimexpr(\linewidth-4\tabcolsep)/2\relax}>{\raggedright\arraybackslash}p{\dimexpr(\linewidth-4\tabcolsep)/2\relax}}
\textbf{A.} 争取民族独立和人民解放是实现民族复兴的前提 & \textbf{B.} 中国已经被卷入世界资本主义经济体系和世界市场中\\[.35em]
\textbf{C.} 中国是一个经济政治发展不平衡的国家 & \textbf{D.} 资本主义制度已经过时\\[.35em]\end{tabular}\end{center}

10.毛泽东在谈到辛亥革命时指出,辛亥革命有它胜利的地方,也有它失败的地方,“辛亥革命把皇帝赶跑,这不是胜利了吗?说它失败,是说辛亥革命只把一个皇帝赶跑。”毛泽东这里所说的“只把一个皇帝赶跑”是指 :


\begin{center}\begin{tabular}{>{\raggedright\arraybackslash}p{\dimexpr(\linewidth-4\tabcolsep)/2\relax}>{\raggedright\arraybackslash}p{\dimexpr(\linewidth-4\tabcolsep)/2\relax}}
\textbf{A.} 没有推翻帝制 & \textbf{B.} 反帝反封建的革命任务没有完成\\[.35em]
\textbf{C.} 孙中山没有继续革命 & \textbf{D.} 袁世凯窃取了胜利果实\\[.35em]\end{tabular}\end{center}

11.1914 年至1918 年的第一次世界大战,是一场空前残酷的大屠杀.它改变了世界政治的格局,也改变了各帝国主义国家在中国的利益格局,对中国产生了巨大的影响。大战使中国的先进分子 :


\begin{center}\begin{tabular}{>{\raggedright\arraybackslash}p{\dimexpr(\linewidth-4\tabcolsep)/2\relax}>{\raggedright\arraybackslash}p{\dimexpr(\linewidth-4\tabcolsep)/2\relax}}
\textbf{A.} 对中国传统文化产生怀疑 & \textbf{B.} 对西方资产阶级民主主义产生怀疑\\[.35em]
\textbf{C.} 认识到工人阶级的重要作用 & \textbf{D.} 认识到必须优先改造国民性\\[.35em]\end{tabular}\end{center}

12.1929 年12 月下旬，红四军党的第九次代表大会在福建上杭县古田村召开，会议总结了红军创立以来的经验，通过了著名的古田会议决议，决议的中心思想是 :


\begin{center}\begin{tabular}{>{\raggedright\arraybackslash}p{\dimexpr(\linewidth-4\tabcolsep)/2\relax}>{\raggedright\arraybackslash}p{\dimexpr(\linewidth-4\tabcolsep)/2\relax}}
\textbf{A.} 中国共产党必须服从共产国际的领导 & \textbf{B.} 武装斗争是中国革命的主要形式\\[.35em]
\textbf{C.} 在农村根据地广泛开展土地革命 & \textbf{D.} 用无产阶级思想进行军队和党的建设\\[.35em]\end{tabular}\end{center}

13.习近平在欧美同学会成立100 周年庆祝大会上的讲话中说：“希望广大留学人员继承和发扬留学报国的光荣传统，做爱国主义的坚守者和传播者，秉持‘先天下之忧而忧，后天下之乐而乐’的人生理想，始终把国家富强、民族振兴、人民幸福作为努力志向，自觉使个人成功的果实结在爱国主义这颗常青树上。”个人成功的果实之所以应该结在爱国主义这棵常青树上，是因为爱国主义是:


\begin{center}\begin{tabular}{>{\raggedright\arraybackslash}p{\dimexpr(\linewidth-4\tabcolsep)/2\relax}>{\raggedright\arraybackslash}p{\dimexpr(\linewidth-4\tabcolsep)/2\relax}}
\textbf{A.} 个人实现人生价值的力量源泉 & \textbf{B.} 人人实现人生价值的直接条件\\[.35em]
\textbf{C.} 个人成功的根本保障 & \textbf{D.} 个人成功的决定性因素\\[.35em]\end{tabular}\end{center}

14.钱学森曾经说过：“我作为一名中国的科技工作者，活着的目的就是为人民服务。如果人民最后对我的一生所做的工作表示满意的话，那才是最高的奖赏。”这说明评价人生价值的根本尺度是 :


\begin{center}\begin{tabular}{>{\raggedright\arraybackslash}p{\dimexpr(\linewidth-4\tabcolsep)/2\relax}>{\raggedright\arraybackslash}p{\dimexpr(\linewidth-4\tabcolsep)/2\relax}}
\textbf{A.} 个体在社会中的地位 & \textbf{B.} 个体在社会中的影响\\[.35em]
\textbf{C.} 个人对社会和他人的生存和发展的贡献 & \textbf{D.} 个体从社会获得的满足尺度\\[.35em]\end{tabular}\end{center}

15.2014 年4 月15 日，中共中央总书记、中央国家安全委员会主席习近平主持召开中央国家安全委员会第一次会议并发表重要讲话，他强调，面对传统安全威胁与非传统安全威胁交织的局面，要准确把握国家安全形式变化新特点新形势，坚持:


\begin{center}\begin{tabular}{>{\raggedright\arraybackslash}p{\dimexpr(\linewidth-8\tabcolsep)/4\relax}>{\raggedright\arraybackslash}p{\dimexpr(\linewidth-8\tabcolsep)/4\relax}>{\raggedright\arraybackslash}p{\dimexpr(\linewidth-8\tabcolsep)/4\relax}>{\raggedright\arraybackslash}p{\dimexpr(\linewidth-8\tabcolsep)/4\relax}}
\textbf{A.} 共同安全观 & \textbf{B.} 亚洲安全观 & \textbf{C.} 总体国家安全观 & \textbf{D.} 地区集体安全观\\[.35em]\end{tabular}\end{center}

16.近年来，中东地区局势持续动荡，恐怖主义、分离主义愈加猖獗，教派矛盾不断升级。尤其是极端恐怖势力于2014 年6 月29 日宣布成立“伊斯兰国”（ISIS），并宣称将建立地跨西亚北非的“哈里发帝国”，对该地区的秩序造成了重大冲击，并且给国家的全球战略带来了影响，这种影响表现为美国 :


\begin{center}\begin{tabular}{>{\raggedright\arraybackslash}p{\dimexpr(\linewidth-4\tabcolsep)/2\relax}>{\raggedright\arraybackslash}p{\dimexpr(\linewidth-4\tabcolsep)/2\relax}}
\textbf{A.} “中东收缩”战略提上议程 & \textbf{B.} “和平演变”战略归于失败\\[.35em]
\textbf{C.} “北约东扩”战略被迫搁置 & \textbf{D.} “重返亚太”战略收到牵制\\[.35em]\end{tabular}\end{center}

\subsection*{二、多项选择题：17～33 题，每小题2 分，共34 分。下列每题给出的四个选项中，至少有两}

个选项是符合题目要求的。请在答题卡上将所选项的字母涂黑。多选或少选均不得分。


17.母质、气候、生物、地形、时间是土壤形成的五大关键成土因素。母质是土壤形成物质基础的初始无机养分的最初来源。气候导致矿物的风化和合成，有机质的形成和积累，土壤中物质的迁移，分解和合成。生物包括植物、动物和微生物等，是促进土壤发生发展最活跃的因素。地形可以使物质在地表进行再分配，使土壤及母质在接受光、热、水等条件方面发生差异。时间是阐明土壤形成发展的历史动态过程，母质、气候、生物和地形等对成工过程的作用随着时间延续而加强。土壤的生产过程说明 :


\begin{center}\begin{tabular}{>{\raggedright\arraybackslash}p{\dimexpr(\linewidth-4\tabcolsep)/2\relax}>{\raggedright\arraybackslash}p{\dimexpr(\linewidth-4\tabcolsep)/2\relax}}
\textbf{A.} 事物总是作为过程而存在 & \textbf{B.} 时间是物质运动的存在形式\\[.35em]
\textbf{C.} 事物的发展总是呈现出线性上升的态势 & \textbf{D.} 事物的产生是多重因素相互作用的结果\\[.35em]\end{tabular}\end{center}

18.平衡是事物发展的一种状态，小到体操中人在平衡木上的行走，杂技中的骑车走钢丝，独轮车表演，直升机在空中的悬停，大到人类的生存，地球的运转，天体的运行等等，都是保持平衡的一种状态，世间的万事万物，之所以能不停的运动、发展、前进，一个重要原因就在于保持了平衡，要使“平衡”成为人们的“大智慧”，就要 :


\begin{center}\begin{tabular}{>{\raggedright\arraybackslash}p{\dimexpr(\linewidth-4\tabcolsep)/2\relax}>{\raggedright\arraybackslash}p{\dimexpr(\linewidth-4\tabcolsep)/2\relax}}
\textbf{A.} 精确把握事物的度 & \textbf{B.} 准确掌握辩证否定的方式和方向\\[.35em]
\textbf{C.} 善于协调事物内部各种因素的相互关系 & \textbf{D.} 全面理解绝对运动和相对静止的辩证关系\\[.35em]\end{tabular}\end{center}

19.马克思主义从必然性和偶然性的辩证统一中理解杰出人物的历史作用,认为 :


\begin{center}\begin{tabular}{>{\raggedright\arraybackslash}p{\dimexpr(\linewidth-4\tabcolsep)/2\relax}>{\raggedright\arraybackslash}p{\dimexpr(\linewidth-4\tabcolsep)/2\relax}}
\textbf{A.} 杰出人物能够改变历史发展的基本方向 & \textbf{B.} 杰出人物的历史作用受到一定历史条件的制约\\[.35em]
\textbf{C.} 杰出人物历史作用的形成和发挥与其服务人民群众的意愿,密不可分 & \textbf{D.} 杰出人物们因其智慧性格因素对社会进程发生影响\\[.35em]\end{tabular}\end{center}

20.马克思指出:“资本主义积累不断地并且同它的能力和规模成比例地生产出相对的,即超过资本增殖的平均需要的,因而是过剩的或追加的工人人口。”“过剩的工人人口是积累或资本主义基础上的财富发展的必然产物,但是这种过剩人口反过来又成为资本主义积累的杠杆,甚至成为资本主义生产方式存在的一个条件.”上述论断表明 :


\begin{center}\begin{tabular}{>{\raggedright\arraybackslash}p{\dimexpr(\linewidth-4\tabcolsep)/2\relax}>{\raggedright\arraybackslash}p{\dimexpr(\linewidth-4\tabcolsep)/2\relax}}
\textbf{A.} 资本主义积累必然导致工人人口的供给相对于资本的需要而过剩 & \textbf{B.} 资本主义社会过剩人口之所以是相对的,是因为它不为资本价值增殖所需要\\[.35em]
\textbf{C.} 资本主义生产周期性特征需要有相对过剩的人口规律与之相适应 & \textbf{D.} 资本主义积累使得资本主义社会的人口失业规模呈现越来越大的趋势\\[.35em]\end{tabular}\end{center}

21.国家垄断资本主义是国家政权和私人垄断资本融合在一起的垄断资本主义。第二次世界大战结束以来,在国家垄断资本主义获得充分发展的同时,资本主义国家通过宏观调节和微观规制对生产、流通、分配和消费各个环节的干预也更加深入。其中,微观规制的类型主要有 :


\begin{center}\begin{tabular}{>{\raggedright\arraybackslash}p{\dimexpr(\linewidth-8\tabcolsep)/4\relax}>{\raggedright\arraybackslash}p{\dimexpr(\linewidth-8\tabcolsep)/4\relax}>{\raggedright\arraybackslash}p{\dimexpr(\linewidth-8\tabcolsep)/4\relax}>{\raggedright\arraybackslash}p{\dimexpr(\linewidth-8\tabcolsep)/4\relax}}
\textbf{A.} 公众生活规制 & \textbf{B.} 公共事业规制 & \textbf{C.} 社会经济规制 & \textbf{D.} 反托拉斯法\\[.35em]\end{tabular}\end{center}

22.2014年5月22日,习近平在上海召开的外国专家座谈会上指出,“要实行更加开放的人才政策,不唯地域引进人才,不求所有开发人才,不拘一格用好人才。”当前,我们之所以比历史上任何时期都更加强调重视人才\textbackslash{}用好人才,是因为人才是 :


\begin{center}\begin{tabular}{>{\raggedright\arraybackslash}p{\dimexpr(\linewidth-4\tabcolsep)/2\relax}>{\raggedright\arraybackslash}p{\dimexpr(\linewidth-4\tabcolsep)/2\relax}}
\textbf{A.} 国家竞争力的重要体现 & \textbf{B.} 先进生产力的集中体现\\[.35em]
\textbf{C.} 第一资源和国家战略资源 & \textbf{D.} 推进中国特色社会主义事业的关键因素\\[.35em]\end{tabular}\end{center}

23.到2012 年底,我国仍在耕地上从事农业生产经营的农民家庭约1.9 亿户，所经营的耕地面积占农村家庭承包耕地总面积的92.5\%；仍实行由集体统一经营的村，组约有2000 个；已发展起农业专业合作社68.9 万个，入社成员5300 多万户；各类农业产业化经营组织30 余万个，带动的农户约1.18 亿户；约有2556 万亩耕地由企业租赁经营，上述现象表明在我国农村 :


\begin{center}\begin{tabular}{>{\raggedright\arraybackslash}p{\dimexpr(\linewidth-4\tabcolsep)/2\relax}>{\raggedright\arraybackslash}p{\dimexpr(\linewidth-4\tabcolsep)/2\relax}}
\textbf{A.} 农业经营主体呈现多样化趋势 & \textbf{B.} 土地的集体所有权性质已经发生变化\\[.35em]
\textbf{C.} 家庭承包经营仍然是最基本的经营形式 & \textbf{D.} 土地经营权的流转使农民失去了对土地的承包权\\[.35em]\end{tabular}\end{center}

24.与十一届全国人民代表大会相比，十二届人代会在代表结构与组成上，呈现“两升一降”的变化，来自一线的工人、农民代表401 名，占代表总数的13.42\%，提高了5.18 个百分点；专业技术人员代表610 名，占代表总数的20.42\%，提高了1.2 个百分点；党政领导干部代表1042名，占代表总数的34.88\%，降低了6.93 个百分点。提高基层人大代表特别是一线工人，农民，知识分子代表比例，降低党政领导干部代表比例，有利于 :


\begin{center}\begin{tabular}{>{\raggedright\arraybackslash}p{\dimexpr(\linewidth-4\tabcolsep)/2\relax}>{\raggedright\arraybackslash}p{\dimexpr(\linewidth-4\tabcolsep)/2\relax}}
\textbf{A.} 推动人民群众最关心最直接最现实问题的解决 & \textbf{B.} 调动基层群众参政议政的积极性与主动性\\[.35em]
\textbf{C.} 保证人民群众直接参加国家管理 & \textbf{D.} 更为充分地发挥全国人大的民意反映与监督职能\\[.35em]\end{tabular}\end{center}

25.国家统计局发布的最新数据显示，2014 年前三季度我国GDP 增长为7.4\%，其中第三季度增长为7.3\%，创下了2009 年一季度以来的新低。总体上看，虽然经济增长有所放缓，但国民经济继续运行在合理区间，稳中有进的态势没有变，今后一个时期经济保持平稳较快发展的可能性仍然较大。这是一种趋势性的变化，是经济到了新的发展阶段表现出来的一种新常态。我国经济新常态的主要特点是:


\begin{center}\begin{tabular}{>{\raggedright\arraybackslash}p{\dimexpr(\linewidth-4\tabcolsep)/2\relax}>{\raggedright\arraybackslash}p{\dimexpr(\linewidth-4\tabcolsep)/2\relax}}
\textbf{A.} 经济増长速度从高速增长转为中高速增长 & \textbf{B.} 中国经济对世界市场的需求减弱\\[.35em]
\textbf{C.} 经济结构不断优化升级 & \textbf{D.} 经济发展动力从要素驱动，投资驱动转向创新驱动\\[.35em]\end{tabular}\end{center}

26.人类社会发展的历史表明，对一个民族，一个国家来说，最持久、最深层次的力量是全社会共同认可的核心价值观。面对世界范围思想文化交流交融交锋形势下价值观的新态势，面对改革开放和发展社会主义市场经济条件下思想意识多元多样多变的新特点，积极培育和践行社会主义核心价值观，有利于 :


\begin{center}\begin{tabular}{>{\raggedright\arraybackslash}p{\dimexpr(\linewidth-4\tabcolsep)/2\relax}>{\raggedright\arraybackslash}p{\dimexpr(\linewidth-4\tabcolsep)/2\relax}}
\textbf{A.} 促进人的全面发展和引领社会全面进步 & \textbf{B.} 巩固全党全国人民团结奋斗的共同思想基础\\[.35em]
\textbf{C.} 巩固马克思主义在意识形态领域的指导地位 & \textbf{D.} 集聚实现中华民族伟大复兴中国梦的强大正能量\\[.35em]\end{tabular}\end{center}

27.甲午，对中国人民和中华民族具有特殊含义，在我国近代史上也具有特殊含义。1894 年7月，日本发动甲午战争，清朝在战争中战败。这场战争对中国的影响主要有：


\begin{center}\begin{tabular}{>{\raggedright\arraybackslash}p{\dimexpr(\linewidth-4\tabcolsep)/2\relax}>{\raggedright\arraybackslash}p{\dimexpr(\linewidth-4\tabcolsep)/2\relax}}
\textbf{A.} 中国海关的行政权落入外国人手中 & \textbf{B.} 中国人开始有了普遍的民族意识觉醒\\[.35em]
\textbf{C.} 台湾被日本侵占 & \textbf{D.} 帝国主义列强掀起瓜分中国的狂潮\\[.35em]\end{tabular}\end{center}

28. 1946 年1 月10 日，政治协商会议在重庆开幕，出席会议的有国民党，共产党，民主同盟，青年党和无党派人士的代表38 人。会议通过了宪法草案，政府组织案，国民大全案，和平建国纲领，军事问题案五项协议。按照协议规定建立的政治体制相当于英国、法国的议会制和内阁制，不是新民主义性质的而且国民党在政府中占着明显的优势。对政协的上述协议，共产党赞同并决心严格履行，这是因为它有利于：


\begin{center}\begin{tabular}{>{\raggedright\arraybackslash}p{\dimexpr(\linewidth-4\tabcolsep)/2\relax}>{\raggedright\arraybackslash}p{\dimexpr(\linewidth-4\tabcolsep)/2\relax}}
\textbf{A.} 推进民主政治的发展和进步 & \textbf{B.} 打破国民党一党独裁的局面\\[.35em]
\textbf{C.} 改变国共两党军事力量对比 & \textbf{D.} 保障解放区政权的合法地位\\[.35em]\end{tabular}\end{center}

29.1979 年,针对当时存在的是否还要坚持毛泽东思想的问题，邓小平指出“有些同志说，我们只拥护‘正确的毛泽东思想’，而不是拥护错误的毛泽东思想。这种说法也是有错误的。”“这种说法”之所以错误，是因为 :


\begin{center}\begin{tabular}{>{\raggedright\arraybackslash}p{\dimexpr(\linewidth-4\tabcolsep)/2\relax}>{\raggedright\arraybackslash}p{\dimexpr(\linewidth-4\tabcolsep)/2\relax}}
\textbf{A.} 没有把毛泽东与党的其他领导人对毛泽东思想的贡献区分开 & \textbf{B.} 没有把毛泽东思想与有中国特色社会主义理论区分开\\[.35em]
\textbf{C.} 没有把毛泽东思想与毛泽东的思想区分开 & \textbf{D.} 没有把毛泽东晚年的错误与毛泽东思想的科学体系区分开\\[.35em]\end{tabular}\end{center}

30.个人品德是通过社会道德教育和个人自觉的道德修养所形成的稳定的心理状态和行为习惯。它是个体对某种道德要求认同和践履的结果，集中体现了道德认知、道德情感、道德意志和道德行为的内在统一，这表明个人品德是:


\begin{center}\begin{tabular}{>{\raggedright\arraybackslash}p{\dimexpr(\linewidth-4\tabcolsep)/2\relax}>{\raggedright\arraybackslash}p{\dimexpr(\linewidth-4\tabcolsep)/2\relax}}
\textbf{A.} 个人行为的统一整体及知、情、意、行的综合体现 & \textbf{B.} 在实践活动中表现出来的行为的稳定倾向\\[.35em]
\textbf{C.} 在实践活动中锻炼而成的一种特殊品性 & \textbf{D.} 偶然的、短暂的道德行为现象\\[.35em]\end{tabular}\end{center}

31.1763 年，老威廉·皮特在《论英国人个人居家安全的权利》的演讲中说：“即使最穷的人，在他的小屋里也能够对抗国王的权威。屋子可能很破旧，屋顶可能摇摇欲坠；风可以吹进这所房子，雨可以淋进这所房子，但是国王不能踏进这所房子，他的千军万马也不敢跨进这件破房子的门槛。”这段话后来被浓缩成“风能进，雨能进，国王不能进”。这凸显了权力与权利的关系是 :


\begin{center}\begin{tabular}{>{\raggedright\arraybackslash}p{\dimexpr(\linewidth-4\tabcolsep)/2\relax}>{\raggedright\arraybackslash}p{\dimexpr(\linewidth-4\tabcolsep)/2\relax}}
\textbf{A.} 权力优先于权利 & \textbf{B.} 权力决定权利\\[.35em]
\textbf{C.} 权力应当以权利为界限 & \textbf{D.} 权力必须受到权利的制约\\[.35em]\end{tabular}\end{center}

32.2014 年2 月27 日，十二届全国人大常委会第七次会议通过决定，将9 月3 日确定为中国人民抗日战争胜利纪念日，将12 月13 日确定为南京大屠杀死难者国家公祭日，设立这两个纪念日：


\begin{center}\begin{tabular}{>{\raggedright\arraybackslash}p{\dimexpr(\linewidth-4\tabcolsep)/2\relax}>{\raggedright\arraybackslash}p{\dimexpr(\linewidth-4\tabcolsep)/2\relax}}
\textbf{A.} 有助于向中国人民和世界各国人民传播历史事实的真相 & \textbf{B.} 是对抗击日本帝国主义侵略付出巨大牺牲和作出巨大贡献的人们的敬重与缅怀\\[.35em]
\textbf{C.} 彰显了中国作为反法西斯主要战场的伟大作用 & \textbf{D.} 是对南京大屠杀中大量死难同胞的告慰和尊重\\[.35em]\end{tabular}\end{center}

33.2014 年11 月5 日至11 日，亚太经济合作组织（APEC）第二十二次领导人非正式会议在北京召开。这是一次开创性的历史盛会，硕果累累，其中，《北京反腐败宣言》的通过尤为引人注目。该《宣言》通过的意义在于，各成员国:


\begin{center}\begin{tabular}{>{\raggedright\arraybackslash}p{\dimexpr(\linewidth-4\tabcolsep)/2\relax}>{\raggedright\arraybackslash}p{\dimexpr(\linewidth-4\tabcolsep)/2\relax}}
\textbf{A.} 加强了涉腐、涉案赃款跨境流动的信息共享 & \textbf{B.} 将形成携手打击跨境腐败的网络\\[.35em]
\textbf{C.} 杜绝了跨国腐败行为的发生 & \textbf{D.} 达成了就追逃、追赃开展执法合作的重要共识\\[.35em]\end{tabular}\end{center}

\subsection*{三、分析题：34～38 小题，每小题10 分，共50 分。要求结合所学知识分析材料并回答问题。}

将答案写在答题纸指定位置上。


34.结合材料回答问题：


\subsection*{材料1}

2014 年11 月19 日，首届世界互联网大会在浙江乌镇召开，习近平总书记在致大会的贺词中指出，互联网真正让世界变成了地球村，让国际社会越来越成为你中有我、我中有你的命运共同体、李克强总理20 日下午在杭州会见出席大会的中外代表并同他们座谈，他表示，互联网是人类最伟大的发明之一，改变了人类世界的空间轴、时间轴和思想维度。中国接入互联网20 年来，已发展成为世界互联网大国，不仅培育起一个巨大市场，也促生了许多新技术、新产品、新业态、新模式，创造了上千万就业岗位，很多人特别是年轻人，大学生因此实现了事业梦，人生梦。

目前，全世界网民数量达到了30 亿人，普及率达40\%，全球范围内实现了网络互联、信息互通。即使是世界上最偏僻的一角，只要接入互联网，就接入了人类这个大家庭。同住地球村的“居民”，借助互联网的力量极大的拉近了距离，互联经济已经成为世界经济发展速度最快、潜力最大、合作最活跃的领域之一，形成了世界网络大市场；一个短小的视频通过全世界网民的点击，可以一夜之间成为全球流行文化的宠儿；提供高速的移动通信和无线宽带服务，几乎已是各国旅游“设施”的标配。

当然，互联网发展过程中也产生了一系类问题，如网络信息安全、网络犯罪等，甚至对国家主权、安全、发展利益形成了新的挑战。互联网到底是阿里巴巴的宝库，还是潘多拉的魔盒？这取决于“命运共同体”如何认真应对，谋求共治。从这样的视野来看，已走过20 年岁月的中国互联网，站在了大有可为的新起点上。摘编自《人民日报》（2014年11 月21 日）

\subsection*{材料2}

中国互联网从1994 年全功能接入国际互联网至今，实现了20 年的高速发展，不仅在技术层面一再突破，而且带来了新的思维理念，有人把他概括为“互联网思维”。对何谓“互联网思维”目前还没有定论。然而，打破思维定势，主动革新自我是互联网思维不变的主题，意味着“便捷、互动、用户至上”等理念。这些理念让人们不断感受到互联网带来的变化与变革。

如果把沃尔玛等传统龙头企业比多大象，那么互联网上的小商户只能叫蚂蚁。数百万只蚂蚁聚合起来，吃掉大象并非没有可能。试想，如果没有互联网，任何一家传统商业企业要想把数百万个商家和近亿客户装进来是不可想象的。

既然去中心的互联网更有利于“蚂蚁”生存，“大象”要彻底摆脱危机，可能就要让自己某种程度上也变成“蚂蚁”，至少自身要具备“蚂蚁”的特性。道理很简单，在互联网环境下与灵活的

“蚂蚁军团”作战，庞大的体量以及传统组织形式带来的大企业病，很可能让“大象”的优势转变为劣势。只有彻底改变基因，让“象”内部产生无数个热衷创业的“蚂蚁”，这仗才有的打，毕竟，和蚂蚁打仗，大象有力用不上，但更庞大、更强大的蚁群则成为最后的胜者。

摘编自《人民日报》（2014 年5 月5 日、5 月26 日等）

（1）联系自身实际，分析为什么“互联网到底是阿里巴巴的宝库，还是潘多拉的魔盒”取决于互联网的“命运共同体”？（5 分）


（2）怎样以辩证的思维方式认识与处理“蚂蚁”与“大象”的关系？（5 分）


35. 结合材料回答问题让大猫小猫都有路走计划经济时期全靠国家管理市场，市场边角被忽略，很多小商品没人去生产，有些新的市场需求也没人去注意。非公有制经济的特点是只要市场有需求，它就会去满足。要保证各种所有制经济依法平等使用生产要素，公平参与市场竞争，同等收到法律保护，就需要重视中小企业融资难的问题，经济学家成思危讲过这样一则寓言：著名科学家牛顿养了两只猫。一只大猫一只小猫。他在墙上开了两个洞，一个大洞。一个小洞，有人笑话他说，你还是大科学家呢，开一个洞就够了，小猫也可以走大洞嘛。牛顿说不对，如果两个猫同时要出去，那大洞肯定被大猫占住了，小猫就无路可走。要真正解决小微企业的问题，就要建立真正为小微企业服务的小型银行，让大银行服务大企业，小银行服务小企业。因为从商业角度说，大银行本身就嫌平爱富，嫌小爱大。小微企业市场风险很大，交易成本也高。跟大企业签一个1 亿元的合同，相当于跟小企业签20 个500 万的合同。现在居民和企业手中有大量存款，而小微企业却贷不到款。这就需要一条通道，这条通道就是社区银行等中小银行。发展民营的、小型的金融机构有利于解决好小微企业的融资困难。大企业是我国经济的脊梁，小微企业是血肉。没有大企业国民经济站不起来，但是如果小微企业垮了，那国民经济不成了骨头架子了吗？十八届三中全会以来，国务院陆续出台了关于大力扶持小微企业健康发展的多项政策，国务院常务会议也多次强调要加快发展民营银行等中小金融机构，为小微企业减负添力。2014 年11 月，李克强总理在浙江考察时再次对民营银行长期致力于服务小微企业给与了充分肯定。为小微企业打开直接融资大门，是开创性的制度安排。一大批有巨大市场潜力的小微企业将会成长为支撑中国经济转型升级的参天大树。

摘编自 光明网(2013 年11 月15 日)、新华网（2014 年11 月21 日）

（1） 现阶段我国发展社会主义市场经济为什么应坚持“让大猫小猫都有路走”？（6 分）


（2） 如何更好地发挥非公有制经济在经济发展中的作用？（4 分）


36.结合材料回答问题：

1925 年郭沫若在一篇文章中讲述了这样一个故事：

十月十五日丁祭过后的第二天，孔子和他的得意门生颜回、子路、子贡三位在上海的文庙里吃着冷肉的时候，有四位年轻的大班抬了一乘朱红漆的四轿，一直闯进庙来，里面走出一位脸如螃蟹，胡须满腮的西洋人来，原来这位胡子螃蟹脸就是马克思。

孔子一见来的是马克思，他便禁不得惊喜着叫出：啊啊，有朋自远方来，不亦乐乎呀！你来到敝庙来，有什么见教呢？

马克思说:我是特为领教而来，我们的主义已经传到你们中国，我希望在你们中国能够实现，但是近来有些人说，我的主义和你的思想不同，所以在你的思想普遍着的中国，我的主义是没有实现的可能性，因此我便来直接领教你：究竟你的思想是怎么样？和我的主义怎样不同？

孔子说：难得你今天亲自到了我这里来，太匆促了，不好请你演讲，至少请你作一番谈话罢。你的理想的世界是怎样的呢？

马克思说：我的理想的世界，是我们生存在这里面，万人要能和一人一样自由平等地发展他们的才能，人人都各能尽力做事而不望报酬，人人都各能得生活的保障而无饥寒的忧虑，这就是我所谓“各尽所能，各取所需”的共产社会。

孔子说：你这个理想社会和我的大同世界竟是不谋而合，你请让我背一段我的旧文章给你听罢。“大道之行也，天下为公，选贤与能，讲信修睦；故人不独亲其亲，不独子其子，使老有所终，壮有所用，幼有所长，矜寡孤独废疾者皆有所养，男有分，女有归；货恶其弃于地也不必藏于己；力恶其不出于身也不必为己；是故谋闭而不兴，盗窃乱贼而不作，故外户而不闭，是谓大同”，这不是和你的理想完全一致的吗？

马克思说我的理想和有些空想家不同，我的理想不是虚构出来的，也并不是一步可以跳到的。我们先从历史上证明社会的产业有逐渐增值之可能，其次是逐渐增值的财产逐渐集中于少数人之手中，于是使社会生出贫乏病来，社会上的争斗便永无宁日。

孔子说：我从前也早就说过“不患寡而患不均，不患贫而患不安”的呀！

孔子的话还没有十分钟落脚，马克思早反对起来了：不对，不对！你和我的见解终竟是两样，我是患寡且患不均，患贫且患不安的，你要晓得，寡了便不均起来，贫了便是不安的根本。所以我对于私产的集中虽是反对，对于产业的增值却不惟不敢反对，而且还极力提倡，所以我们一方面用莫大的力量去剥夺私人的财产，而同时也要以莫大的力量来增值社会的产业。孔子说：尊重物质本是我们中国的传统思想，洪范八政食货为先，管子也说过“仓源实而知礼节，衣食足而知荣辱”，我的思想乃至我国的传统思想，根本和你一样，总要先把产业提高起来，然后才来均分。

马克思到此才感叹起来：我不想在两千年前，在遥远的东方，已经有了你这样的一个老同志！你我的见解完全一致的，怎么有人曾说我的思想和你的不合，和你们中国的国情不合，不

能施行于中国呢？


（1）“马克思进文庙”的历史背景是什么？（4 分）


（2）如何理解孔子与马克思对话中谈到的他们之间思想上的“不同”与“一致”？（6 分）


37. 结合材料回答问题：


\subsection*{材料1}

2014 年10 月闭幕的十八届四中全会，是党在中央全会上第一次专题讨论已发治国的问题，体现了对法治的高度重视。会议结束后，微博上的各种评论满是对法治进步的热望：“想要法治的果实，就要给它阳光雨露”“期待法治进入与人民互动的2.0 时代”“法治不仅是宏大的，也是具体的；它关乎国家治理，更关于百姓福祉”……《韩非子》有句名言“国无常强，无常弱。奉法者强则国强，奉法者弱则过弱”。尊奉法律，需要执政者，治理者发力，引导之，提倡之，遵守法律，需要全体公民给力，用法律来定分止争，维护之，践行之。网络上已经有人以普通人“小明”为例，演绎“四中全会与你我有啥关系”。有认识，法治于人就如同空气，你可能不会时时刻刻意识到它的存在，可一旦缺少就立刻窒息。的确，从出生到成长，从成家到立业，舞步需要法治的护航：加强对财产权的保护，完善教育、医疗、视频安全等方面的法律法规，提高环境污染的违法成本……四中全会促动“法治的春天”有着温暖人心的春意。当越来越多人在法治的护佑下感受着畅快的呼吸，法治才能成为内心时时恪守的律令。

也不用回避中国的法治还有很多问题，从“暂行50 多年”的高温条例，到保护个人信息安全等方面尚无完善法律，中国的法治进程需要紧跟时代的步伐。四中全会从立法、司法、执法、守法等方面开出来药方，但最根本的还是提升全社会对法治的信心与信任，正如党的十八届四中全会公报所说，法律的权威源自于人民内心的拥护和真诚信仰。这才是法治的力量所在，尊严所系。摘自《人民日报》2014 年10 月24 日

\subsection*{材料2}

法治是人类为了征服自己，由人类自己立法进行自我管理，这远比征服自然困难得多。特别是约束公权力，非有高度的觉悟，顽强的毅力和坚强的意志难以成其事。任何国家法治的确立都不是在一盘散沙的状态下随随便便建立起来的，而是必须有坚定有力的集中统一领导和部署。迄今为止，尚未有法治成功的国家是在群龙无首，四分五裂的状态下实现法治的。恰恰相反，就法治发达国家的经验来看，这些国家的法治之所以能够最终确立，都是自上而下，从官到民表现出对法治制作的追求，付出巨大的努力。在中国这个拥有13 亿人口，情况极其复杂的大国建设法治，更需要有自上而下将强统一的领导，要有统一的意志，坚决果断一体推行。正是基于这样的情况，十八界四中全会指出，全面推进依法治国，必须坚持党的领导。


摘自《人民日报》2014 年10 月29 日


（1）如何理解“法治关乎国家治理，更关乎百姓福祉”？（6 分）


（2）为什么“全面推进依法治国，必须坚持党的领导”（4 分）


38. 阅读下列材料


\subsection*{材料1}

1960 年5 月27 日，毛泽东与来华访问的英国元帅蒙哥马利，围绕“50 年以后中国的命运”有一段深刻的对话。蒙哥马利说，我有一个有趣的问题想问下主席：中国大概需要50 年，一切事情就办得差不多了。到那时候，你看中国的前途将会怎样？历史的教训是，当一个国家非常强大的时候，就倾向于侵略，是不是？要向外国侵略，就会被打回来；外国是外国人住的地方，别人不能去，没有权利也没有理由硬挤进去，如果去，就要被赶走，这是历史教训。如果我们占人家一寸土地，我们就是侵略者。

“蒙哥马利之问”折射的是一些西方人内心深处的“国强必霸”逻辑。然而，这样的逻辑与中国人千百年来的民族心理完全不在一个“频道”上。正如习近平所说，“中华民族的血液中没有侵略他人，称霸世界的基因。”600 多年前，郑和受命出使西洋，足迹遍布30 多个国家和地区。明朝初期的中国，是综合国力位居世界前列的强国。但是，与地理大发现时期欧洲国家的殖民政策不同。郑和船队始终奉行“共享太平之福”的宗旨，尊重当地习惯，平等开展多边贸易，把中国的建筑、绘画、雕刻、服饰等领域的精湛技术带入亚非国家，促进了中外文化的双向交流和共同进步。郑和下西洋的“和平之旅”永载史册。

明朝洪武年间，缅甸与百夷(今缅甸北部)交战，明太祖未发一兵，派李思聪、钱古训二人劝和。二人先奉劝缅甸“两国之民居处虽分，惟存关市之讥。是其和也，其或纷争不已，天将昭鉴福善祸淫”，又告诫百夷“莫如守全，以图绵长，不亦美乎”。双方均为道义所感悟，最终罢战息兵。

孙中山先生曾在日本演说，“东方的文化是王道，主张仁义道德，西方的文化是霸道，主张功利强权。讲仁义道德，是由正义公理来感化人，讲功利强权，是用洋枪大炮来压迫人。”近代中国遭受列强欺凌，无数仁人志士高喊“落后就要挨打”“振兴中华”，但只是为了获得免于被欺凌的自由，为了以平等的姿态屹立于世界民族之林。正如“和平学之父”约翰·加尔通所说，有些人总希望有一个暴力选择，但中国以自己特有的视角来观察现实，阴阳平衡、尊重智慧、众生平等理念被视为理所当然，和平关系的普遍原则以相互合作、平等互利为起点。

拿破仑说，中国是一头沉睡的狮子，当这头睡狮醒来时，世界都会为之发抖。今年3 月，习近平同志在法国巴黎向世界宣示，中国这头狮子已经醒了，但这是一只和平的、可亲的、文明的狮子。读懂了“和”文化是中国人千百年来流淌的血脉，就感受到了走向世界的中国那种无

法改变的“和”的气度与内质。

摘编自《人民日报》（2010 年12 月22 日，2014 年5 月20 日）

（1）毛泽东和蒙哥马利的“对话”反映了什么？（4 分）


（2）如何理解习近平所说“中国这头狮子已经醒了，但这是一只和平的、可亲的、文明的狮子”？（6 分）

\end{document}
